\subsection{Entropy and the Arrow of Time}
  \label{subsec:entropy-arrow-of-time}

  The arrow of time is a fundamental structural feature arising from the
  irreversibility of the reprojection sequence
  $\chi_{\mathrm{eff}}(t) \to \chi_{\mathrm{eff}}(t+1)$, not a derived statistical phenomenon.
  Each reprojection step loses information under the non-injective
  projection~$\Pi$, making exact reversal structurally inadmissible:
  no admissible successor state can restore the full pre-image of its predecessor.
  $\chi$ does not itself relax; it supplies the admissibility constraints
  that orient the sequence of constrained reconstructions.
  Irreversibility arises from information loss under~$\Pi$; directionality
  arises from the structural constraints of~$\chi$.
  This is consistent with the monotone decrease of the relational
  energy $E_\chi$ along admissible transitions established in
  Section~\ref{sec:temporal-succession}: that monotonicity is the projected image of the
  structural ordering imposed by~$\chi$, not a property of $\chi$ itself evolving in time.

  Entropy increase emerges only at the level of effective spacetime
  descriptions, providing a statistical summary of how macroscopic degrees
  of freedom evolve under successive non-injective reprojections.
  Entropy growth does not explain the arrow of time; it reflects the
  cumulative information loss accumulated under the reprojection sequence.
  The temporal asymmetry is not present in $\chi$ itself---which is
  atemporally fixed---but in the structure of admissible reprojections that $\chi$ constrains.

  This reverses the standard explanatory hierarchy: time asymmetry is
  imposed structurally by the admissibility constraints on the reprojection
  sequence, not attributed to special initial conditions or to a substrate that itself evolves.
  Processes involving deprojection do not correspond to entropy decrease
  or temporal reversal but represent a transition to a level of description
  where thermodynamic notions no longer apply.
